This section fixes the single canonical object used in the main proof spine.

\subsection*{Classical definitions (standard)}

\begin{definition}[Critical strip and critical line]
The critical strip is
\[
\{s\in\C:0<\RePart(s)<1\},
\]
and the critical line is
\[
\{s\in\C:\RePart(s)=\tfrac12\}.
\]
\end{definition}

\begin{definition}[Reflection involution]
Let $\sigma:\C\to\C$ be $\sigma(s)=1-s$.
\end{definition}

\subsection*{Canonical manuscript object}

\begin{definition}[Main variable and classical slice]
The main variable in this manuscript is $q\in\C$.
When we compare with classical notation, we use the complex slice identification $q=s$.
\end{definition}

\begin{definition}[Canonical balance object]
\label{def:phi-canonical}
Define
\[
\Phi(q):=\Gamma(q)(1-q)\zeta(q).
\]
This is the only canonical main-text object used for balance.
\end{definition}

\begin{definition}[Balanced point]
\label{def:balanced-point}
A point $q$ is \emph{balanced} if and only if
\[
\Phi(q)=0.
\]
Inside the critical strip, write
\[
\mathcal{B}:=\{q\in\C:0<\RePart(q)<1,\ \Phi(q)=0\}.
\]
\end{definition}

\begin{remark}[Immediate classical correspondence on the strip]
For $0<\RePart(q)<1$, the gamma factor has no zeros and $1-q\neq0$.
Hence on the strip,
\[
\Phi(q)=0 \iff \zeta(q)=0.
\]
The manuscript uses this equivalence as the direct mathematical reason for choosing $\Phi$.
\end{remark}

\begin{remark}[Status boundary]
The choice
\[
\Phi(q)=\Gamma(q)(1-q)\zeta(q)
\]
is an intentional simplification for audit.
It is a modeling choice in this manuscript under review, not a proof that this is the only possible formulation.
No RH conclusion is claimed at this definitional stage.
\end{remark}
